\documentclass[11pt]{article}

\usepackage[margin=1in]{geometry}
\usepackage[T1]{fontenc}
\usepackage{lmodern}
\usepackage{amsmath,amssymb,amsthm}
\usepackage{mathtools}
\usepackage{enumitem}
\usepackage{array}
\usepackage{microtype}
\usepackage{tikz}
\usepackage[hidelinks]{hyperref}

\newcommand{\R}{\mathbb{R}}
\newcommand{\one}{\mathbf{1}}
\newcommand{\hull}{\operatorname{hull}}
\newcommand{\lexle}{\le_{\mathrm{lex}}}
\newcommand{\lexlt}{<_{\mathrm{lex}}}
\DeclareMathOperator{\conv}{conv}
\DeclareMathOperator{\relint}{relint}
\DeclareMathOperator{\aff}{aff}
\DeclareMathOperator{\cl}{cl}
\DeclareMathOperator{\inter}{int}
\DeclareMathOperator{\spt}{spt}

\theoremstyle{plain}
\newtheorem{theorem}{Theorem}[section]
\newtheorem{proposition}[theorem]{Proposition}
\newtheorem{lemma}[theorem]{Lemma}
\newtheorem{corollary}[theorem]{Corollary}
\theoremstyle{definition}
\newtheorem{definition}[theorem]{Definition}
\newtheorem{example}[theorem]{Example}

\setlist{itemsep=2pt,topsep=4pt}

\title{Coherence and accuracy for sets of desirable gambles}
\author{Jason Konek and Catrin Campbell-Moore}
\date{October 2026}
\hypersetup{pdftitle={Coherence and accuracy for sets of desirable gambles},
  pdfauthor={Jason Konek and Catrin Campbell-Moore}}

\begin{document}
\maketitle

\section{Setting and definitions}\label{sec:setting}

\subsection{Worlds, gambles and states}

\begin{definition}[Worlds and gambles]\label{def:gambles}
Let $\Omega$ be a nonempty set, of any cardinality; its elements are \emph{worlds}. Let
$\ell^\infty(\Omega)$ be the space of bounded real functions on $\Omega$, with the norm
$\|h\|\coloneqq\sup_{\omega\in\Omega}|h(\omega)|$ and the pointwise order: $h\ge g$ means
$h(\omega)\ge g(\omega)$ for every $\omega\in\Omega$. We write $g\gneq0$ if $g\ge0$ and $g\ne0$.
Write $\one$ for the constant function
with value $1$. A \emph{gamble space} is a norm-closed, separable linear subspace
$E\subseteq\ell^\infty(\Omega)$ that contains $\one$. Throughout, $E$ is a fixed gamble space,
and its elements are \emph{gambles}; the gamble $h$ pays $h(\omega)$ at the world $\omega$. For
$h\in E$ and $r>0$, $B(h,r)\coloneqq\{g\in E:\|g-h\|<r\}$ is the open ball; $\cl D$ and
$\inter D$ are the norm closure and the norm interior of $D\subseteq E$. A \emph{convex
cone} is a set $D\subseteq E$ with $D+D\subseteq D$ and $\lambda D\subseteq D$ for every
$\lambda>0$. For $D,D'\subseteq E$ and $c>0$ we write $D+D'\coloneqq\{d+d':d\in D,\ d'\in D'\}$,
$h+D\coloneqq\{h+d:d\in D\}$, $h-D\coloneqq\{h-d:d\in D\}$, $-D\coloneqq\{-d:d\in D\}$ and
$D/c\coloneqq\{d/c:d\in D\}$.
\end{definition}

\begin{definition}[Coherent sets of desirable gambles]\label{def:coherentsets}
A set $D\subseteq E$ is a \emph{coherent set of desirable gambles} if
\begin{itemize}
  \item[(D1)] $0\notin D$;
  \item[(D2)] $g\in D$ for every $g\in E$ with $g\gneq0$;
  \item[(D3)] $\lambda g\in D$ whenever $g\in D$ and $\lambda>0$;
  \item[(D4)] $g+h\in D$ whenever $g,h\in D$.
\end{itemize}
Thus a coherent set of desirable gambles is a convex cone that excludes $0$ and contains every
nonzero gamble that is nowhere negative.
\end{definition}

\begin{definition}[States and credal sets]\label{def:states}
A \emph{state} is a linear functional $q:E\to\R$ that is positive ($q(h)\ge0$ whenever
$h\ge0$) and normalised ($q(\one)=1$). We write $S$ for the set of states, with the
weak-$\ast$ topology, the coarsest topology for which $q\mapsto q(h)$ is continuous for every
$h\in E$; $E^\ast$ is the space of norm-continuous linear functionals on $E$, and $\conv$
denotes the convex hull. For each world $\omega$, the evaluation $\delta_\omega(h)\coloneqq h(\omega)$ is a
state. A \emph{credal set} is a nonempty, weak-$\ast$ closed, convex set $Q\subseteq S$; $S$
itself is a credal set. For a credal set $Q$ and $h\in E$ put
\[
  P_Q(h)\coloneqq\min_{q\in Q}q(h),\qquad D_Q^\circ\coloneqq\{h\in E:P_Q(h)>0\},\qquad
  C_Q\coloneqq\{h\in E:P_Q(h)\ge0\},
\]
\[
  D_Q\coloneqq D_Q^\circ\cup\{g\in E:g\gneq0\};
\]
the minimum exists by Lemma~\ref{lem:states}(a). For a state $q$ we write $\{q>0\}\coloneqq
\{h\in E:q(h)>0\}$, and similarly $\{q<0\}$ and $\{q=0\}$; thus $\{q>0\}=D^\circ_{\{q\}}$.
\end{definition}

\begin{lemma}\label{lem:states}
Let $Q,Q_1,Q_2$ be credal sets.
\begin{enumerate}[label=(\alph*)]
  \item Every state has norm $1$. $S$ is convex, weak-$\ast$ compact and metrizable; so every
  credal set is weak-$\ast$ compact, and the minimum defining $P_Q$ is attained.
  \item For every $h\in E$, $\min_{q\in S}q(h)=\inf_\Omega h$ and $\max_{q\in S}q(h)=
  \sup_\Omega h$. Hence $D_S^\circ=\{h:\inf_\Omega h>0\}$, $-D_S^\circ=\{h:\sup_\Omega h<0\}$
  and $D_S=\{h:h\gneq0\}$.
  \item $P_Q$ is $1$-Lipschitz, superadditive, positively homogeneous and monotone ($P_Q(h)\ge
  P_Q(g)$ whenever $h\ge g$), and $P_Q(h+t\one)=P_Q(h)+t$ for $t\in\R$. Hence $D_Q^\circ$ is an
  open convex cone, $C_Q$ is closed, $D_S^\circ\subseteq D_Q^\circ$, and neither $D_Q^\circ$
  nor $D_Q$ contains an $h$ with $\sup_\Omega h\le0$. Moreover $\conv(Q_1\cup Q_2)$ is a credal
  set, $D^\circ_{Q_1}\cap D^\circ_{Q_2}=D^\circ_{\conv(Q_1\cup Q_2)}$ and $D_{Q_1}\cap D_{Q_2}=
  D_{\conv(Q_1\cup Q_2)}$.
  \item For $0\le\tau\le1$, $(1-\tau)Q+\tau S\coloneqq\{(1-\tau)q+\tau s:q\in Q,\ s\in S\}$ is a
  credal set, and $P_{(1-\tau)Q+\tau S}=
  (1-\tau)P_Q+\tau P_S$.
  \item There is a countable set $\Omega_0\subseteq\Omega$ with $\inf_{\Omega_0}h=\inf_\Omega h$
  for every $h\in E$.
\end{enumerate}
\end{lemma}

\begin{lemma}\label{lem:DQ}
Let $Q$ be a credal set. Then $D_Q$ is a coherent set of desirable gambles and a Borel set,
$D_Q^\circ$ is the interior of $D_Q$, and $D_Q^\circ\subseteq D_Q\subseteq C_Q$ with $D_Q\setminus
D_Q^\circ\subseteq\{P_Q=0\}$.
\end{lemma}

\begin{proof}
If $g\gneq0$, then $P_Q(g)\ge P_Q(0)=0$ by monotonicity (Lemma~\ref{lem:states}(c)). Hence
$D_Q\subseteq C_Q$, and every $h\in D_Q\setminus D_Q^\circ$ has $h\gneq0$ and $P_Q(h)=0$.
(D1): $P_Q(0)=0$, and $0\gneq0$ fails. (D2) holds by definition. (D3): $P_Q$ is positively
homogeneous, and $\lambda g\gneq0$ whenever $g\gneq0$ and $\lambda>0$. (D4): let $g,h\in D_Q$.
If $g,h\gneq0$, then $g+h\ge g\gneq0$, so $g+h\gneq0$. Otherwise one of them, say $g$, lies in
$D_Q^\circ$, and $P_Q(h)\ge0$; by superadditivity, $P_Q(g+h)\ge P_Q(g)+P_Q(h)>0$. The set
$\{g:g\ge0\}=\bigcap_{\omega\in\Omega}\{g:g(\omega)\ge0\}$ is closed, so $\{g:g\gneq0\}$ is
Borel; $D_Q^\circ$ is open (Lemma~\ref{lem:states}(c)); so $D_Q$ is Borel. Finally, $D_Q^\circ$
is an open subset of $D_Q$. If $h\in D_Q\setminus D_Q^\circ$ and $t>0$, then $P_Q(h-t\one)=-t
<0$, so $h-t\one\notin D_Q^\circ$, and $h-t\one\not\ge0$, since otherwise $P_Q(h-t\one)\ge0$; so
$h-t\one\notin D_Q$. As $\|(h-t\one)-h\|=t$, $h$ is not an interior point of $D_Q$.
\end{proof}

\subsection{The measure}

\begin{definition}[The measure]\label{def:measure}
Throughout, $\mu$ is a finite Borel measure on $(E,\|\cdot\|)$, and we use the same letter for
its completion. A set is \emph{null} if it is measurable for the completion of $\mu$ and has
measure $0$; a statement holds \emph{almost everywhere} (a.e.) if it holds outside a null set.
Each gamble $f$ represents a binary choice, and a model judges $f$ and $-f$ separately
(Definition~\ref{def:advice}); each choice therefore offers two opportunities to give a
judgement, one for each way round. The measure $\mu$ gives the relative frequency or importance
of these opportunities: $\mu(B)$ is the weight of the opportunities to judge the gambles in $B$.
We consider the following conditions on $\mu$.
\begin{itemize}
  \item[(M1)] \emph{Finite first moment:} $\int\|h\|\,d\mu(h)<\infty$.
  \item[(M2)] \emph{Null indifference boundaries:} $\mu\{P_Q=0\}=0$ for every credal set $Q$.
  \item[(M3)] \emph{Full support:} every nonempty open subset of $E$ has positive measure.
  \item[(M4)] \emph{Reflection:} for every Borel set $B\subseteq E$, $\mu(B)=0$ if and only if
  $\mu(-B)=0$.
\end{itemize}
Condition (M4) says that whether a choice opportunity is negligible should not depend on which
way round it is written: $f$ and $-f$ represent the same choice (Definition~\ref{def:advice}).
Since $B\mapsto-B$ is an involution, (M4) holds as soon as $\mu(B)=0$ implies $\mu(-B)=0$ for
every Borel $B$. Every \emph{symmetric} measure, one with $\mu(-B)=\mu(B)$ for every Borel $B$,
satisfies (M4). Under (M4), $-X$ is null whenever $X$ is null.
\end{definition}

\begin{definition}[Atoms and diffuse part]\label{def:atoms}
An \emph{atom} of $\mu$ is a gamble $h$ with $\mu\{h\}>0$. Let $A_\mu$ be the set of atoms. The
\emph{diffuse part} of $\mu$ is $\mu_c(B)\coloneqq\mu(B\setminus A_\mu)$, for every set $B$
that is measurable for the completion of $\mu$; this is well defined because $A_\mu$ is
countable (Lemma~\ref{lem:tools}(a)). The measure $\mu$ is \emph{atomless} if
$A_\mu=\emptyset$, that is, if $\mu_c=\mu$. We also consider the condition
\begin{itemize}
  \item[(M3$'$)] \emph{Full support of the diffuse part:} every nonempty open subset of $E$ has
  positive $\mu_c$-measure.
\end{itemize}
Condition (M3$'$) implies (M3); for atomless $\mu$ the two coincide.

Each result below holds uniformly over every $\mu$ that satisfies the conditions the result
states; beyond those conditions, no particular distribution of advice opportunities is
assumed. The main results are not claimed for every finite Borel measure:
Theorem~\ref{thm:dominance} and Corollary~\ref{cor:admissible} assume (M2) and (M3$'$);
Proposition~\ref{prop:axioms} assumes (M1); the last claim of Corollary~\ref{cor:existence}
and the results of Section~\ref{sec:repair} assume (M1) and (M2). These
conditions are substantive assumptions. They say the following.
\begin{center}
\small
\begin{tabular}{@{}>{\raggedright\arraybackslash}p{0.07\textwidth}>{\raggedright\arraybackslash}p{0.5\textwidth}>{\raggedright\arraybackslash}p{0.35\textwidth}@{}}
\hline
& \textbf{What it requires} & \textbf{Used by}\\
\hline
(M1) & The stakes of the advice opportunities, $\|h\|$, have finite average size. & The score and
the integral ranking (Section~\ref{sec:integral}) and Section~\ref{sec:repair}; not
Theorem~\ref{thm:dominance} or Corollary~\ref{cor:admissible}.\\[2pt]
(M2) & For each credal set $Q$, the gambles on the boundary of $D_Q^\circ$, those with
$P_Q(h)=0$, carry no weight. & Theorem~\ref{thm:dominance}, Corollary~\ref{cor:admissible},
the last claim of Corollary~\ref{cor:existence}, and Section~\ref{sec:repair}, where by
Lemma~\ref{lem:identify}(b) it makes the coherent models exactly the models equal a.e.\ to
some $D_Q$; not Proposition~\ref{prop:axioms}.\\[2pt]
(M3) & No region of gambles is negligible: every nonempty open set carries positive weight. &
No result assumes it directly; it follows from (M3$'$), and coincides with (M3$'$) for an
atomless measure, such as that of Lemma~\ref{lem:measure}.\\[2pt]
(M3$'$) & Every nonempty open set carries positive weight even after the weight on single
gambles (the atoms) is discarded, so that opportunities are spread through every region
rather than concentrated at isolated points. & Theorem~\ref{thm:dominance} and
Corollary~\ref{cor:admissible}; no result in Sections~\ref{sec:integral}
or~\ref{sec:repair} assumes it.\\[2pt]
(M4) & Whether a set of choice opportunities is negligible does not depend on which way round
the choices are written. & Only Proposition~\ref{prop:advice}, which justifies reading models as
advice; no other result uses it.\\
\hline
\end{tabular}
\end{center}
\end{definition}


\begin{lemma}\label{lem:measure}
For every gamble space $E$ there is an atomless, symmetric finite Borel measure on $E$ that
satisfies (M1), (M2) and (M3), and hence (M3$'$) and (M4).
\end{lemma}

\subsection{Models and coherence}

\begin{definition}[Models and coherence]\label{def:models}
A \emph{model} is a set $X\subseteq E$ that is measurable for the completion of $\mu$; it
\emph{accepts} the gambles it contains. Two sets $X,Y\subseteq E$ are \emph{equal a.e.} if their
symmetric difference $X\triangle Y$ is null; then $X$ is a model if and only if $Y$ is. Models
that are equal a.e.\ are identified. We write
$X\ne Y$ a.e.\ if $X$ and $Y$ are not equal a.e. A model is \emph{non-null} if it is not null. A model $X$ is
\emph{coherent} if $X=D$ a.e.\ for some coherent set of desirable gambles $D$
(Definition~\ref{def:coherentsets}), and \emph{incoherent} otherwise. By Lemma~\ref{lem:DQ},
every model equal a.e.\ to some $D_Q$ is coherent; under (M2), every coherent model is of this
form (Lemma~\ref{lem:identify}(b)).
\end{definition}

\subsection{Models as advice}

\begin{definition}[Advice]\label{def:advice}
A gamble $f$ represents the binary choice between any two gambles $g$ and $h$ with $f=g-h$:
at the world $\omega$, taking $g$ rather than $h$ gains $f(\omega)$, and taking $h$ rather
than $g$ gains $-f(\omega)$. The gambles $f$ and $-f$ represent the same choice, written the
two ways round. By accepting $f$, a model judges $f$ desirable, that is, judges taking $g$
rather than $h$ worthwhile; by accepting $-f$, it judges the reverse. Its advice on the
choice is fixed by both judgements. For a model $X$ and a gamble $f$, the \emph{advice} of
$X$ on $f$, written $\alpha_X(f)$, is one of four \emph{advice states}:
\begin{enumerate}[label=(\roman*)]
  \item \emph{strict advice for $f$}, if $f\in X$ and $-f\notin X$;
  \item \emph{strict advice against $f$}, if $-f\in X$ and $f\notin X$;
  \item \emph{silence}, if $f\notin X$ and $-f\notin X$;
  \item \emph{contradictory advice}, if $f\in X$ and $-f\in X$.
\end{enumerate}
Thus $\alpha_X(-f)$ is $\alpha_X(f)$ with ``for'' and ``against'' exchanged. A coherent set of
desirable gambles $D$ never gives contradictory advice, that is, never contains both $f$ and
$-f$: if $f\in D$ and $-f\in D$, then $0=f+(-f)\in D$ by (D4), contradicting (D1). Under (M4), a coherent model gives contradictory
advice only on a null set of choices (Proposition~\ref{prop:advice}).
\end{definition}

\begin{proposition}\label{prop:advice}
Write $\check\mu(B)\coloneqq\mu(-B)$. The following are equivalent.
\begin{enumerate}[label=(\alph*)]
  \item $\mu$ satisfies (M4).
  \item For all models $X,Y$ with $X=Y$ a.e., the set $\{f:\alpha_X(f)\ne\alpha_Y(f)\}$ is null.
  \item For all models $X,Y$ with $X=Y$ a.e., the set $\{f:\alpha_X(f)\ne\alpha_Y(f)\}$ is
  $(\mu+\check\mu)$-null.
  \item There is a $\sigma$-ideal $\mathcal N$ of subsets of $E$ such that, for all models $X,Y$,
  $X=Y$ a.e.\ if and only if $\{f:\alpha_X(f)\ne\alpha_Y(f)\}\in\mathcal N$.
\end{enumerate}
Under (M4), $X\mapsto\alpha_X$ is a bijection, up to null sets, from models onto the
$\mu$-measurable maps $\alpha$ from $E$ to the four advice states such that, for a.e.\ $f$,
$\alpha(-f)$ is $\alpha(f)$ with ``for'' and ``against'' exchanged.
\end{proposition}

\begin{proof}
For models $X,Y$ put $N\coloneqq X\triangle Y$. The advice of $X$ and $Y$ on $f$ differs exactly
when $f\in N$ or $-f\in N$, so
\[
  \{f:\alpha_X(f)\ne\alpha_Y(f)\}=N\cup(-N). \tag{$\ast$}
\]
(a)$\Rightarrow$(c). If $N$ is null, then $N\subseteq N_0$ for a Borel set $N_0$ with
$\mu(N_0)=0$. By (M4), $M\coloneqq N_0\cup(-N_0)$ has $\mu(M)=0$; as $M=-M$, also
$\check\mu(M)=0$. So $N\cup(-N)\subseteq M$ is $(\mu+\check\mu)$-null.
(c)$\Rightarrow$(b) since $\mu\le\mu+\check\mu$.
(b)$\Rightarrow$(d). Let $\mathcal N$ be the null sets. One direction is (b); conversely, if
$N\cup(-N)$ is null then so is $N$.
(d)$\Rightarrow$(a). Let $B$ be Borel with $\mu(B)=0$. With $X\coloneqq B$ and $Y\coloneqq
\emptyset$, ($\ast$) and (d) give $B\cup(-B)\in\mathcal N$. With $X\coloneqq-B$, a Borel set,
and $Y\coloneqq\emptyset$, the same set appears in ($\ast$), so (d) gives $-B=\emptyset$
a.e., that is, $\mu(-B)=0$.

Bijection. Assume (M4). By (b), $\alpha_X$ depends, up to null sets, only on the class of
$X$. Also, $-X$ is a model whenever $X$ is: if $X=B\cup M$ with $B$ Borel and $M$
null, then $-X=(-B)\cup(-M)$ with $-M$ null. So $\alpha_X$ is measurable. If $\alpha_X=\alpha_Y$
a.e., then $X\triangle Y\subseteq\{\alpha_X\ne\alpha_Y\}$ is null. Given $\alpha$ as stated,
let $X$ be the set of $f$ with $\alpha(f)$ strict advice for $f$ or contradictory advice, and
let $Z$ be the null set of $f$ at which $\alpha(-f)$ is not $\alpha(f)$ with ``for'' and
``against'' exchanged. For $f\notin Z$, $-f\in X$ if and only if $\alpha(f)$ is strict advice
against $f$ or contradictory advice; so $\alpha_X(f)=\alpha(f)$ for $f\notin Z$.
\end{proof}

\begin{definition}[Unordered choices]\label{def:unordered}
A set $B\subseteq E$ with $B=-B$ is \emph{symmetric}; it represents a set of unordered choices.
The symmetric Borel sets form a $\sigma$-algebra $\mathcal S$. A \emph{distribution over
unordered choices} is a finite measure $\pi$ on $\mathcal S$. Its \emph{equal split} is
\[
  \mu_\pi(B)\coloneqq\int\tfrac12\bigl(\one_B(f)+\one_B(-f)\bigr)\,d\pi(f)\qquad(B\text{ Borel}),
\]
which gives each way round of a choice half the weight of the choice.
\end{definition}

\begin{lemma}\label{lem:unordered}
Let $\pi$ be a distribution over unordered choices.
\begin{enumerate}[label=(\alph*)]
  \item $\mu_\pi$ is the only symmetric Borel measure that agrees with $\pi$ on $\mathcal S$, and
  $\mu_\pi(B)=\tfrac12\pi(B\cup(-B))+\tfrac12\pi(B\cap(-B))$ for every Borel $B$.
  \item Let $\mu$ be a Borel measure that agrees with $\pi$ on $\mathcal S$. Then
  $(\mu+\check\mu)/2=\mu_\pi$, and $\mu$ has the same null sets as $\mu_\pi$ if and only if $\mu$
  satisfies (M4).
  \item If $\mu$ is symmetric and satisfies (M1), then the score of Definition~\ref{def:score}
  satisfies $m_\omega(X)=\tfrac12\int\bigl(\one_X(f)-\one_X(-f)\bigr)f(\omega)\,d\mu(f)$ for
  every model $X$ and world $\omega$.
\end{enumerate}
\end{lemma}

\begin{proof}
For Borel $B$, $f\mapsto\one_B(f)+\one_B(-f)$ equals $\one_{B\cup(-B)}+\one_{B\cap(-B)}$, so it
is $\mathcal S$-measurable. Hence $\mu_\pi$ is a Borel measure (by monotone convergence), it is
symmetric, it agrees with $\pi$ on $\mathcal S$, and the formula in (a) holds. If $\mu$ agrees
with $\pi$ on $\mathcal S$, then, as the integrand is $\mathcal S$-measurable,
$\tfrac12(\mu(B)+\mu(-B))=\int\tfrac12(\one_B(f)+\one_B(-f))\,d\mu(f)=\mu_\pi(B)$; this is the
first claim of (b), and for symmetric $\mu$ it gives $\mu=\mu_\pi$, the uniqueness in (a). By
(a), $\mu_\pi(B)\le\pi(B\cup(-B))=\mu(B\cup(-B))\le2\mu_\pi(B)$, so $\mu_\pi(B)=0$ if and only
if $\mu(B\cup(-B))=0$. Under (M4), $\mu(B)=0$ if and only if $\mu(B\cup(-B))=0$; conversely,
the null sets of the symmetric measure $\mu_\pi$ are closed under reflection. (c) For symmetric
$\mu$, $\int g(f)\,d\mu(f)=\int g(-f)\,d\mu(f)$ for integrable $g$; apply this to
$g(f)\coloneqq\one_X(f)f(\omega)$ and average.
\end{proof}

\begin{example}\label{ex:noreflection}
Conditions (M1)--(M3$'$) do not imply (M4). Let $\mu_0$ be the measure of
Lemma~\ref{lem:measure} and $\mu\coloneqq\mu_0+\delta_{\one}$, where $\delta_{\one}$ is the unit
mass at $\one$. Then $\mu$ satisfies (M1), (M2) and (M3$'$), and $\{-\one\}$ is null while
$\{\one\}$ is not. The coherent set of desirable gambles $D_S$ and the model $D_S\cup\{-\one\}$
are equal a.e., so the second is a coherent model; yet the first gives strict advice for $\one$
and the second gives contradictory advice on the same choice, which has positive measure.
\end{example}

\begin{definition}[Support]\label{def:spt}
The \emph{support} of a model $X$ is
\[
  \spt X\coloneqq\{h\in E:\mu_c(X\cap B(h,\varepsilon))>0\text{ for every }\varepsilon>0\}.
\]
\end{definition}

\begin{lemma}\label{lem:tools}
\begin{enumerate}[label=(\alph*)]
  \item The set $A_\mu$ is countable, and $\mu_c$ is a finite Borel measure with $\mu_c\{h\}=0$
  for every $h\in E$.
  \item If $B$ is a model and $0\le t\le\mu_c(B)$, there is a Borel set $G\subseteq B\setminus
  A_\mu$ with $\mu(G)=t$.
  \item For every model $X$, $\spt X$ is closed and $\mu_c(X\setminus\spt X)=0$. If $X\subseteq
  Y$ a.e., then $\spt X\subseteq\spt Y$.
  \item Assume (M2). Then $A_\mu\subseteq D_S^\circ\cup(-D_S^\circ)$.
\end{enumerate}
\end{lemma}

\begin{lemma}\label{lem:identify}
\begin{enumerate}[label=(\alph*)]
  \item Let $D\subseteq E$ be a convex cone with $0\notin D$ and $D_S^\circ\subseteq D$. Then
  $K_D\coloneqq\{q\in S:q(h)\ge0\text{ for all }h\in D\}$ is a credal set, and
  $D^\circ_{K_D}\subseteq D\subseteq C_{K_D}$.
  \item Let $D$ be a coherent set of desirable gambles. Then $K_D$ is a credal set and
  $D_{K_D}\subseteq D\subseteq C_{K_D}$. Under (M2), $D$ is a model and $D=D_{K_D}$ a.e.; so,
  under (M2), a model is coherent if and only if it is equal a.e.\ to $D_Q$ for some credal set
  $Q$.
\end{enumerate}
\end{lemma}

\subsection{Supporting states and the coherent hull}

\begin{definition}[Supporting states and coherent hull]\label{def:hull}
For a model $B$, the \emph{supporting states} of $B$ are
\[
  Q_B\coloneqq\{q\in S:q(h)\ge0\text{ for a.e.\ }h\in B\}.
\]
The model $B$ is \emph{consistent} if $Q_B\ne\emptyset$ and \emph{inconsistent} otherwise. The
\emph{coherent hull} of $B$ is $\hull(B)\coloneqq D_{Q_B}$ if $B$ is consistent, and
$\hull(B)\coloneqq E$ if $B$ is inconsistent.
\end{definition}

\begin{lemma}\label{lem:hull}
Let $B,B'$ be models.
\begin{enumerate}[label=(\alph*)]
  \item $Q_B$ is weak-$\ast$ closed and convex, so it is a credal set if $B$ is consistent. If
  $B\subseteq B'$ a.e., then $Q_{B'}\subseteq Q_B$.
  \item If $B$ is consistent, then $P_{Q_B}\ge0$ a.e.\ on $B$. Under (M2), $B\subseteq\hull(B)$
  a.e.
  \item Let $Q$ be a credal set. If $B\subseteq D_Q$ a.e., then $Q\subseteq Q_B$ and
  $\hull(B)\subseteq D_Q$.
\end{enumerate}
\end{lemma}

\subsection{Changes}

\begin{definition}[Changes]\label{def:change}
Let $X$ and $Y$ be models. The \emph{change from $Y$ to $X$} adds the gambles
$A\coloneqq X\setminus Y$ and removes the gambles $B\coloneqq Y\setminus X$; it is \emph{null}
if $A\cup B$ is null, that is, if $X=Y$ a.e.
\begin{itemize}
  \item The change is \emph{available} to a model $Z$, and $Z$ \emph{could make} it, if
  $B\subseteq Z$ a.e.\ and $A\cap Z$ is null.
  \item The change is \emph{guided at} the world $\omega$ if $A\subseteq\{h:h(\omega)>0\}$
  a.e.\ and $B\subseteq\{h:h(\omega)<0\}$ a.e. It is \emph{guided at every world} if it is
  guided at $\omega$ for every $\omega\in\Omega$.
\end{itemize}
\end{definition}

\subsection{Margin triples and consolidation}

\begin{definition}[Margin triples]\label{def:triple}
Let $c\in\{1,2\}$. A \emph{margin triple of order $c$} is a triple $(G_1,G_2,G_3)$ of pairwise
disjoint models such that
\begin{enumerate}[label=(\roman*)]
  \item $\mu(G_1)=\mu(G_2)>0$ and $\mu(G_3)=c\,\mu(G_1)$; and
  \item there is $\kappa>0$ such that $z(\omega)\ge\bigl(x(\omega)+y(\omega)\bigr)/c+\kappa$ for
  every $\omega\in\Omega$, every $x\in G_1$, every $y\in G_2$ and every $z\in G_3$.
\end{enumerate}
Condition (ii) refers to the sets $G_1,G_2,G_3$ themselves, not only to their classes up to
null sets. A margin triple of order $1$ is a \emph{sum triple}, and one of order $2$ is a
\emph{midpoint triple}.
\end{definition}

\begin{definition}[Consolidation of a model]\label{def:consolidation}
Let $(G_1,G_2,G_3)$ be a margin triple, and let $W$ be a model with $G_1\cup G_2\subseteq W$
a.e.\ and $W\cap G_3$ null. The \emph{consolidation} of $W$ by $(G_1,G_2,G_3)$ is the model
\[
  W'\coloneqq\bigl(W\setminus(G_1\cup G_2)\bigr)\cup G_3 .
\]
It depends only on $W$ up to null sets. The change from $W$ to $W'$ adds $G_3$ and removes
$G_1\cup G_2$, up to null sets.
\end{definition}

\begin{lemma}\label{lem:Wincoherent}
Let $(G_1,G_2,G_3)$ be a margin triple of order $c$, and let $\kappa$ be as in
Definition~\ref{def:triple}(ii). No coherent model $W$ satisfies $G_1\cup G_2\subseteq W$
a.e.\ and $W\cap G_3$ null.
\end{lemma}

\begin{proof}
Suppose that a coherent model $W$ does, and let $D$ be a coherent set of desirable gambles with
$W=D$ a.e. Since $G_1\setminus D\subseteq(G_1\setminus W)\cup(W\setminus D)$ is null and $G_1$
is non-null, $G_1\cap D$ is non-null; so there is $x\in G_1\cap D$. Likewise there is $y\in
G_2\cap D$. By (D4) and (D3), $(x+y)/c\in D$. For every $z\in G_3$, the gamble $g\coloneqq
z-(x+y)/c$ satisfies $g\ge\kappa\one$, so $g\gneq0$ and $g\in D$ by (D2); hence $z=g+(x+y)/c\in
D$ by (D4). So $G_3\subseteq D$. Then $G_3\setminus W\subseteq D\setminus W$ is null, and as
$G_3$ is non-null, $G_3\cap W$ is not null, a contradiction.
\end{proof}

\section{Rankings and dominance}\label{sec:rankings}

\begin{definition}[Rankings]\label{def:ranking}
A \emph{ranking} assigns to each world $\omega$ a binary relation $\succeq_\omega$ on models.
Each instance $X\succeq_\omega Y$ is a \emph{comparison} at $\omega$, read ``$X$ is at least as
good as $Y$ at $\omega$''. The \emph{strict comparison} $X\succ_\omega Y$ means that
$X\succeq_\omega Y$ and not $Y\succeq_\omega X$. The models $X$ and $Y$ are \emph{tied} at
$\omega$, written $X\sim_\omega Y$, if $X\succeq_\omega Y$ and $Y\succeq_\omega X$.
\end{definition}

\begin{definition}[Dominance]\label{def:dominance}
Fix a ranking, and let $X$ and $Y$ be models.
\begin{itemize}
  \item $Y$ \emph{dominates} $X$ if $Y\succeq_\omega X$ at every world $\omega$ and
  $Y\succ_\omega X$ at some world $\omega$.
  \item $Y$ \emph{strictly dominates} $X$ if $Y\succ_\omega X$ at every world $\omega$.
  \item $X$ is \emph{undominated}, or \emph{admissible}, if no model dominates $X$.
\end{itemize}
\end{definition}

\section{Axioms}\label{sec:axioms}

\begin{definition}[Axioms]\label{def:axioms}
Let $\mathcal K$ be a family of coherent models. A ranking may satisfy the following axioms.
In each, $X,Y,Z,W,A,B,C$ range over models and $\omega$ over worlds.
\begin{itemize}[leftmargin=3.2em]
  \item[(A1)] \emph{Guidance.} If $X\ne Y$ a.e.\ and the change from $Y$ to $X$ is guided at
  every world, then $X\succ_\omega Y$ for every $\omega$.

  \emph{Gloss.} A change that is not null, adds only gambles that pay off at every world, and
  removes only gambles that lose at every world makes a model strictly better at every world.

  \item[(A2)] \emph{Three-model condition.} There are no $\omega$, $A$, $B$, $C$ with
  $A\succeq_\omega B$, $B\succeq_\omega C$ and $C\succ_\omega A$.

  \emph{Gloss.} If $A$ is at least as good as $B$, and $B$ at least as good as $C$, at a
  world, then $C$ is not strictly better than $A$ at that world.

  \item[(A3)] \emph{Transitivity through ties.} If $Z\sim_\omega X$ and $Y\succeq_\omega Z$,
  then $Y\succeq_\omega X$.

  \emph{Gloss.} A model that is at least as good as $Z$ at a world is at least as good, at
  that world, as every model tied with $Z$ there.

  \item[(A4)] \emph{Consolidation.} Let $(G_1,G_2,G_3)$ be a margin triple of order $1$ or
  $2$, and let $W$ be a model with $G_1\cup G_2\subseteq W$ a.e.\ and $W\cap G_3$ null. Then
  the consolidation $W'$ of $W$ by $(G_1,G_2,G_3)$ satisfies $W'\succeq_\omega W$ for every
  $\omega$.

  \emph{Gloss.} Suppose that a model accepts the gambles in $G_1$ and $G_2$, accepts none of those in
  $G_3$, and each gamble in $G_3$ pays, at every world, at least a fixed amount more than the
  sum (order $1$) or the average (order $2$) of any gamble in $G_1$ and any gamble in $G_2$;
  then giving up $G_1$ and $G_2$ and accepting $G_3$ instead does not make the model worse at
  any world.

  \emph{Explanation.} Read through Definitions~\ref{def:advice} and~\ref{def:measure}, (A4) is a principle about the value of the advice that a model
  supplies. Write $m\coloneqq\mu(G_1)$ and $\mu|_G(B)\coloneqq\mu(B\cap G)$.
  \begin{itemize}[leftmargin=1.2em]
    \item \emph{The two forms.} A sum triple has masses $m,m,m$: judgements of weight $m$ in
    each of $G_1$ and $G_2$ are withdrawn, and judgements of weight $m$ in $G_3$ are supplied;
    at every world, the gamble in each supplied judgement pays at least $\kappa$ more than the
    sum of the gambles in any withdrawn pair. A midpoint triple has masses $m,m,2m$: twice that
    weight is supplied, and the gamble in each supplied judgement pays at least $\kappa$ more
    than the average of the gambles in any withdrawn pair.
    \item \emph{What the masses supply.} Equal masses fix how much of each part is replaced.
    They do not by themselves fix how often choices from $G_1$ and $G_2$ occur together. The
    comparison is made precise by a \emph{grouping of opportunities for comparison}: a finite
    measure on $G_1\times G_2\times G_3$ (sum triples) or on $G_1\times G_2\times G_3\times
    G_3$ (midpoint triples) whose marginals on the first two factors are $\mu|_{G_1}$ and
    $\mu|_{G_2}$ and whose marginals on the remaining factors add up to $\mu|_{G_3}$.
    For a sum triple, $\mu|_{G_1}\otimes\mu|_{G_2}\otimes\mu|_{G_3}/m^2$ on $G_1\times G_2\times
    G_3$ has total mass $m$ and marginals exactly $\mu|_{G_1}$, $\mu|_{G_2}$ and $\mu|_{G_3}$;
    in each group $(x,y,z)$, $z(\omega)\ge x(\omega)+y(\omega)+\kappa$ at every world. The
    midpoint form matches two original opportunities with two replacement ones:
    $\mu|_{G_1}\otimes\mu|_{G_2}\otimes\mu|_{G_3}\otimes\mu|_{G_3}/(4m^3)$ on $G_1\times
    G_2\times G_3\times G_3$ has total mass $m$, marginals $\mu|_{G_1}$ and $\mu|_{G_2}$ on the
    first two factors, and $\tfrac12\mu|_{G_3}$ on each of the last two, which together give
    $\mu|_{G_3}$; in each group $(x,y,z,z')$, $z(\omega)+z'(\omega)\ge x(\omega)+y(\omega)+
    2\kappa$ at every world. Because the payoff inequality holds for every group, the particular
    grouping does not matter.
    \item \emph{Every context.} The axiom applies whatever else $W$ accepts, including where
    adding or removing a gamble creates or resolves contradictory advice. The axiom compares the
    advice supplied before and after; it does not
    treat silence as a default choice of either option.
    \item \emph{Judgements, not whole choices.} The grouped units are single judgements, each
    the acceptance of one gamble, weighted by $\mu$; a choice offers two of them, one for each
    way round (Definition~\ref{def:measure}). When some $G_i$ meets $-G_j$ (possibly with
    $i=j$), both judgements on one choice fall in the parts, in opposite orientations. For
    instance, if $-z\in G_1$ and $z\in G_3$, the consolidation reverses the advice on that
    choice. The payoff inequality still compares single judgements, so these cases need no
    separate treatment. Two of them have direct consequences. If $z$ and $-z$ both lie in
    $G_3$, then $x+y\le-c\kappa\one$ for all $x\in G_1$ and $y\in G_2$: contradictory advice is
    supplied only in place of withdrawn pairs that lose at least $c\kappa$ at every world. If
    $x\in G_1$ and $-x\in G_2$, so that contradictory advice is withdrawn, then $z\ge\kappa\one$
    for every $z\in G_3$.
    \item \emph{The remaining commitment.} Payoff advantages can compensate across advice
    opportunities: the advantage of the supplied judgements at some choices can outweigh what
    the withdrawn judgements secured at others. The measure makes the quantities comparable; it
    does not establish that compensation, which is the substantive content of (A4). By
    Lemma~\ref{lem:Wincoherent}, (A4) constrains the comparison only for incoherent $W$.
  \end{itemize}

  \item[(A5)] \emph{Admissibility of $\mathcal K$.} Every member of $\mathcal K$ is
  undominated.

  \emph{Gloss.} No model is at least as good as a member of $\mathcal K$ at every world and
  strictly better than it at some world.
\end{itemize}
\end{definition}

\section{The dominance theorem}\label{sec:dominance}

\begin{theorem}\label{thm:dominance}
Assume (M2) and (M3$'$), and let a ranking satisfy (A1)--(A4). Then every incoherent model is
dominated.
\end{theorem}

The proof is in Section~\ref{sec:proof}.

\subsection{Ties}

\begin{lemma}\label{lem:tie}
Assume (A2) and (A3). If $Z\sim_\omega X$ at every world $\omega$ and $Y$ dominates $Z$, then
$Y$ dominates $X$.
\end{lemma}

\begin{proof}
By (A3), $Y\succeq_\omega X$ at every $\omega$. Let $\omega_0$ be a world with
$Y\succ_{\omega_0}Z$. If $X\succeq_{\omega_0}Y$, then $Z\succeq_{\omega_0}X$,
$X\succeq_{\omega_0}Y$ and $Y\succ_{\omega_0}Z$, which (A2) excludes. So $Y\succ_{\omega_0}X$.
\end{proof}

\begin{lemma}\label{lem:constie}
Assume (A2)--(A4). Let $X$ be undominated, and let $X'$ be the consolidation of $X$ by a margin
triple of order $1$ or $2$. Then $X'\sim_\omega X$ at every world $\omega$, and $X'$ is
undominated.
\end{lemma}

\begin{proof}
By (A4), $X'\succeq_\omega X$ at every $\omega$. If $X\succeq_\omega X'$ failed at some
$\omega$, then $X'\succ_\omega X$ there, and $X'$ would dominate $X$. So $X'\sim_\omega X$ at
every $\omega$. If some $Y$ dominated $X'$, then Lemma~\ref{lem:tie}, applied with $Z\coloneqq
X'$, would make $Y$ dominate $X$.
\end{proof}

\subsection{Guidance}

\begin{definition}[Sure-gain conditions]\label{def:U}
A model $X$ satisfies
\begin{itemize}
  \item[(U1)] if $D_S^\circ\setminus X$ and $X\cap(-D_S^\circ)$ are null;
  \item[(U2)] if there are no $x\in\spt X$ and $z\in\spt(E\setminus X)$ with $z-x\in D_S^\circ$.
\end{itemize}
\end{definition}

\begin{lemma}\label{lem:U1}
Assume (A1). Every undominated model satisfies (U1).
\end{lemma}

\begin{proof}
Let $X$ be a model and $Y\coloneqq(X\cup D_S^\circ)\setminus(-D_S^\circ)$. Since $D_S^\circ\cap
(-D_S^\circ)=\emptyset$, the change from $X$ to $Y$ adds $D_S^\circ\setminus X$ and removes
$X\cap(-D_S^\circ)$. Every added gamble
is positive at every world, and every removed gamble is negative at every world
(Lemma~\ref{lem:states}(b)); so the change is guided at every world. If $X$ violates (U1), the
change is not null, and (A1) gives $Y\succ_\omega X$ at every $\omega$; so $Y$ dominates $X$.
\end{proof}

\begin{lemma}\label{lem:U2}
Assume (M3$'$). Let $X$ satisfy (U1) but not (U2). Then there is a sum triple $(G_1,G_2,G_3)$
with $G_1\cup G_2\subseteq X$, $G_3\cap X=\emptyset$ and $G_2\subseteq D_S^\circ$, and the
consolidation of $X$ by it violates (U1).
\end{lemma}

\begin{proof}
Take $x\in\spt X$ and $z\in\spt(E\setminus X)$ with $z-x\in D_S^\circ$. Put $4\delta\coloneqq
\inf_\Omega(z-x)>0$, $d\coloneqq(z-x)/2$ and $\rho\coloneqq\delta/2$. For $x'\in B(x,\rho)$,
$y'\in B(d,\rho)$ and $z'\in B(z,\rho)$,
\[
  \inf_\Omega(z'-x'-y')\ge\inf_\Omega(z-x-d)-3\rho=\inf_\Omega d-3\rho=2\delta-\tfrac32\delta=
  \tfrac12\delta .
\]
Also $\inf_\Omega y'\ge2\delta-\rho>0$, so $B(d,\rho)\subseteq D_S^\circ$. Since $x\in\spt X$,
$\mu_c(X\cap B(x,\rho))>0$. By (M3$'$), $\mu_c(B(d,\rho))>0$, and by (U1),
$\mu_c(X\cap B(d,\rho))=\mu_c(B(d,\rho))$. Since $z\in\spt(E\setminus X)$,
$\mu_c(B(z,\rho)\setminus X)>0$. Choose
\[
  0<m<\min\bigl\{\mu_c(X\cap B(x,\rho)),\ \tfrac12\mu_c(B(d,\rho)),\ \mu_c(B(z,\rho)\setminus
  X)\bigr\}.
\]
By Lemma~\ref{lem:tools}(b) there are Borel sets $G_1\subseteq(X\cap B(x,\rho))\setminus
A_\mu$, then $G_2\subseteq(X\cap B(d,\rho))\setminus(G_1\cup A_\mu)$ (which has
$\mu_c$-measure greater than $m$), and $G_3\subseteq B(z,\rho)\setminus(X\cup A_\mu)$, each
with $\mu$-measure $m$. They are pairwise disjoint, since $G_1\cap G_2=\emptyset$ and $G_3\cap
X=\emptyset$. With $\kappa\coloneqq\delta/2$, $(G_1,G_2,G_3)$ is a sum triple, and $G_2\subseteq
D_S^\circ$. The consolidation $(X\setminus(G_1\cup G_2))\cup G_3$ misses $G_2$, so it misses a subset
of $D_S^\circ$ of measure $m>0$, and it violates (U1).
\end{proof}

\begin{corollary}\label{cor:U}
Assume (M3$'$) and (A1)--(A4). Every undominated model satisfies (U1) and (U2).
\end{corollary}

\begin{proof}
Let $X$ be undominated. By Lemma~\ref{lem:U1}, $X$ satisfies (U1). If $X$ violated (U2),
Lemma~\ref{lem:U2} would give a consolidation of $X$ that violates (U1); it is undominated by
Lemma~\ref{lem:constie}, which contradicts Lemma~\ref{lem:U1}.
\end{proof}

\subsection{Identification of coherent models}

\begin{lemma}\label{lem:open}
Assume (M3$'$). Let $X$ satisfy (U1) and (U2), and put $V\coloneqq\spt X+D_S^\circ$.
\begin{enumerate}[label=(\alph*)]
  \item $V$ is open and $\mu_c(V\setminus X)=0$.
  \item $\mu_c(X\setminus\cl V)=0$.
  \item $\cl V+D_S^\circ\subseteq V$.
  \item $D_S^\circ\subseteq V$, and $V$ contains no $h$ with $\sup_\Omega h\le0$.
  \item Assume (M2) in addition. If $V+V\subseteq V$ and $(V+V)/2\subseteq V$, then $X$ is
  coherent.
\end{enumerate}
\end{lemma}

\begin{proof}
(a) $V=\bigcup_{x\in\spt X}(x+D_S^\circ)$ is a union of open sets. Suppose $\mu_c(V\setminus X)>0$,
and put $A\coloneqq V\setminus X$. By Lemma~\ref{lem:tools}(c), $\mu_c(A\setminus\spt A)=0$, so
some $z\in A$ lies in $\spt A$. Since $A\subseteq E\setminus X$, $z\in\spt(E\setminus X)$.
Since $z\in V$, $z-x\in D_S^\circ$ for some $x\in\spt X$. This contradicts (U2).

(b) By Lemma~\ref{lem:tools}(c), $\mu_c(X\setminus\spt X)=0$. For $x\in\spt X$ and $t>0$,
$x+t\one\in V$, and $x+t\one\to x$ as $t\to0$; so $\spt X\subseteq\cl V$.

(c) Since $D_S^\circ+D_S^\circ\subseteq D_S^\circ$, $V+D_S^\circ\subseteq V$. Let $v\in\cl V$
and $d\in D_S^\circ$, and take $v_n\in V$ with $v_n\to v$. Since $D_S^\circ$ is open,
$v-v_n+d\in D_S^\circ$ for large $n$, and then $v+d=v_n+(v-v_n+d)\in V+D_S^\circ\subseteq V$.

(d) Let $d\in D_S^\circ$. For small $\varepsilon>0$, $B(d,\varepsilon)\subseteq D_S^\circ$; then
$\mu_c(B(d,\varepsilon))>0$ by (M3$'$), and $\mu_c(B(d,\varepsilon)\setminus X)=0$ by (U1). So
$d\in\spt X$. Hence $D_S^\circ\subseteq\spt X$. For $d\in D_S^\circ$ and small $t>0$,
$d-t\one\in D_S^\circ\subseteq\spt X$, so $d=(d-t\one)+t\one\in V$. Now suppose $v\in V$ and
$\sup_\Omega v\le0$. As $V$ is open, $v-t\one\in V$ for some $t>0$, and $\sup_\Omega(v-t\one)
<0$; so $V\cap(-D_S^\circ)$ is a nonempty open set, and $\mu_c(V\cap(-D_S^\circ))>0$ by (M3$'$). But
$\mu_c(V\setminus X)=0$ by (a), and $X\cap(-D_S^\circ)$ is null by (U1). This is a contradiction.

(e) \emph{$V$ is convex.} Let $v,w\in V$. Closure under midpoints gives $(1-\lambda)v+\lambda w
\in V$ for every dyadic rational $\lambda\in(0,1)$, by induction on the denominator. Let
$\lambda\in(0,1)$ and $p\coloneqq(1-\lambda)v+\lambda w$. Choose $r>0$ with $B(v,r)\subseteq
V$, and a dyadic rational $\lambda'\in(0,1)$ with $|\lambda-\lambda'|\,\|w-v\|<r(1-\lambda')$.
Put $v'\coloneqq(p-\lambda'w)/(1-\lambda')$. Then $v'-v=(\lambda-\lambda')(w-v)/(1-\lambda')$,
so $v'\in B(v,r)\subseteq V$, and $p=(1-\lambda')v'+\lambda'w\in V$.

\emph{$V$ is a convex cone.} Let $v\in V$. For $0<\lambda<1$, choose $\eta>0$ with
$v-\eta\one\in V$ and put $\varepsilon\coloneqq\lambda\eta/(1-\lambda)$; then $\varepsilon\one
\in D_S^\circ\subseteq V$ by (d), and $\lambda v=\lambda(v-\eta\one)+(1-\lambda)\varepsilon\one\in V$
by convexity. Closure under sums gives $2^kv\in V$ for every $k\ge1$; writing $\lambda\ge1$ as
$2^k\lambda''$ with $0<\lambda''\le1$ gives $\lambda v\in V$. So $V$ is a convex cone.

\emph{Conclusion.} By (d), $0\notin V$ and $D_S^\circ\subseteq V$. By
Lemma~\ref{lem:identify}(a), $K\coloneqq K_V$ is a credal set with $D_K^\circ\subseteq V\subseteq
C_K$. Since $C_K$ is closed, $\cl V\setminus V\subseteq C_K\setminus D_K^\circ=\{P_K=0\}$, which
is null by (M2); also $V\setminus D_K^\circ\subseteq\{P_K=0\}$. With (a) and (b),
\[
  \mu_c(X\triangle D_K^\circ)\le\mu_c(V\setminus X)+\mu_c(X\setminus\cl V)+\mu(\cl V\setminus V)+
  \mu(V\setminus D_K^\circ)=0 .
\]
It remains to treat the atoms. By Lemma~\ref{lem:tools}(d), each atom $a$ lies in $D_S^\circ$ or
in $-D_S^\circ$. If $a\in D_S^\circ$, then $a\in D_K^\circ$ (Lemma~\ref{lem:states}(c)) and $a\in
X$, since $D_S^\circ\setminus X$ is null by (U1). If $a\in-D_S^\circ$, then $a\notin D_K^\circ$
(Lemma~\ref{lem:states}(c)) and $a\notin X$, since $X\cap(-D_S^\circ)$ is null. So $X\triangle
D_K^\circ$ contains no atom, and $\mu(X\triangle D_K^\circ)=\mu_c(X\triangle D_K^\circ)=0$. By
Lemma~\ref{lem:DQ}, $D_K\setminus D_K^\circ\subseteq\{P_K=0\}$, which is null by (M2); so $X=D_K$
a.e. As $D_K$ is a coherent set of desirable gambles (Lemma~\ref{lem:DQ}), $X$ is coherent.
\end{proof}

\begin{lemma}\label{lem:gap}
Assume (M3$'$). Let $X$ satisfy (U1) and (U2), put $V\coloneqq\spt X+D_S^\circ$, let $c\in\{1,2\}$,
and let $v_1,v_2\in V$ with $s\coloneqq(v_1+v_2)/c\notin V$. Then there is a margin triple
$(G_1,G_2,G_3)$ of order $c$ with $G_1\cup G_2\subseteq X$ and $G_3\cap X=\emptyset$, and the
consolidation of $X$ by it violates (U2).
\end{lemma}

\begin{proof}
\emph{Parameters.} Choose $\eta>0$ with $B(v_i,2\eta)\subseteq V$ for $i=1,2$, and put
$r\coloneqq\eta/8$, $v_i'\coloneqq v_i-\eta\one$ and $z_0\coloneqq s-(\eta/c)\one$.

\emph{Three balls.} For $i=1,2$, $B(v_i',r)\subseteq B(v_i,\eta+r)\subseteq V$, so
$\mu_c(B(v_i',r)\setminus X)=0$ by Lemma~\ref{lem:open}(a). For $z\in B(z_0,r)$,
$\inf_\Omega(s-z)\ge\eta/c-r>0$, so $z\in s-D_S^\circ$. The set $s-D_S^\circ$ misses $\cl V$: if $w\in\cl V$
and $s-w\in D_S^\circ$, then $s\in\cl V+D_S^\circ\subseteq V$ by Lemma~\ref{lem:open}(c). Hence
$B(z_0,r)\cap\cl V=\emptyset$, and $\mu_c(B(z_0,r)\cap X)=0$ by Lemma~\ref{lem:open}(b).

\emph{Margin.} Since $z_0-(v_1'+v_2')/c=(\eta/c)\one$, every $x\in B(v_1',r)$, $y\in
B(v_2',r)$ and $z\in B(z_0,r)$ satisfy
\[
  \inf_\Omega\Bigl(z-\frac{x+y}c\Bigr)\ge\frac\eta c-r-\frac{2r}c=
  \begin{cases}5\eta/8,&c=1,\\ \eta/4,&c=2.\end{cases}
\]

\emph{A ball below $B(v_1',r)$.} Put $J_0\coloneqq B(v_1'-3r\one,r)$. Then $J_0\subseteq
B(v_1,\eta+4r)\subseteq V$, so $\mu_c(J_0\setminus X)=0$. For $x\in B(v_1',r)$ and $w\in J_0$,
$\inf_\Omega(x-w)>3r-2r=r$; so $x-w\in D_S^\circ$, and $J_0\cap B(v_1',r)=\emptyset$.

\emph{The triple.} By (M3$'$) the balls $B(v_1',r)$, $B(v_2',r)$, $B(z_0,r)$ and $J_0$ have
positive $\mu_c$-measure. Choose
\[
  0<m<\min\bigl\{\mu_c(B(v_1',r)),\ \tfrac12\mu_c(B(v_2',r)),\ \tfrac1c\mu_c(B(z_0,r)),\
  \tfrac12\mu_c(J_0)\bigr\}.
\]
By Lemma~\ref{lem:tools}(b) there are Borel sets $G_1\subseteq(X\cap B(v_1',r))\setminus
A_\mu$ and $G_2\subseteq(X\cap B(v_2',r))\setminus(G_1\cup A_\mu)$ of $\mu$-measure $m$, and
$G_3\subseteq B(z_0,r)\setminus(X\cup A_\mu)$ of $\mu$-measure $cm$. They are pairwise disjoint,
since $G_1\cap G_2=\emptyset$ and $G_3\cap X=\emptyset$; with $\kappa\coloneqq\eta/4$,
$(G_1,G_2,G_3)$ is a margin triple of order $c$.

\emph{The consolidation violates (U2).} Let $X'\coloneqq(X\setminus(G_1\cup G_2))\cup G_3$ and
$J\coloneqq(J_0\cap X)\setminus G_2$. Then $\mu_c(J)\ge\mu_c(J_0)-m>0$. The set $J$ misses $G_1$
(as $J_0\cap B(v_1',r)=\emptyset$) and $G_2$, and lies in $X$; so $J\subseteq X'$. By
Lemma~\ref{lem:tools}(c) some $w\in J$ lies in $\spt J\subseteq\spt X'$, and some $x\in G_1$
lies in $\spt G_1$. Since $G_1\subseteq X$ and $G_1\cap(G_2\cup G_3)=\emptyset$, $G_1\cap
X'=\emptyset$, so $\spt G_1\subseteq\spt(E\setminus X')$. Finally $x\in B(v_1',r)$ and $w\in
J_0$, so $x-w\in D_S^\circ$. Hence $X'$ violates (U2).
\end{proof}

\subsection{Proof of Theorem~\ref{thm:dominance}}\label{sec:proof}

\begin{proof}[Proof of Theorem~\ref{thm:dominance}]
Suppose that $X$ is incoherent and undominated. By Corollary~\ref{cor:U}, $X$ satisfies (U1)
and (U2). Put $V\coloneqq\spt X+D_S^\circ$. Since $X$ is incoherent, Lemma~\ref{lem:open}(e) shows that
$V+V\not\subseteq V$ or $(V+V)/2\not\subseteq V$: there are $c\in\{1,2\}$ and $v_1,v_2\in V$
with $(v_1+v_2)/c\notin V$. Lemma~\ref{lem:gap} gives a consolidation $X'$ of $X$, by a margin
triple of order $c$, that violates (U2). By Lemma~\ref{lem:constie}, $X'$ is undominated, so it
satisfies (U2) by Corollary~\ref{cor:U}. This is a contradiction.
\end{proof}

\begin{corollary}\label{cor:admissible}
Assume (M2) and (M3$'$), and let a ranking satisfy (A1)--(A4).
\begin{enumerate}[label=(\alph*)]
  \item If the ranking also satisfies (A5) for a family $\mathcal K$ of coherent models, then
  every member of $\mathcal K$ is undominated, while every incoherent model is dominated.
  \item In particular, if the ranking also satisfies (A5) for the family of all coherent
  models, then a model is undominated if and only if it is coherent.
\end{enumerate}
\end{corollary}

\begin{proof}
(a) The first claim is (A5), and the second is Theorem~\ref{thm:dominance}. (b) Coherent
models are undominated by (A5), and incoherent models are dominated by
Theorem~\ref{thm:dominance}.
\end{proof}

\section{The integral ranking}\label{sec:integral}

\begin{definition}[Score]\label{def:score}
Assume (M1). Since $E$ is separable, the identity map of $E$ is Bochner integrable with respect
to $\mu$ \cite{DiestelUhl}. For a model $X$ put
\[
  m(X)\coloneqq\int_Xh\,d\mu(h)\in E,\qquad m_\omega(X)\coloneqq m(X)(\omega)=
  \int_Xh(\omega)\,d\mu(h)\quad(\omega\in\Omega).
\]
The number $m_\omega(X)$ is the \emph{score} of $X$ at $\omega$. Bounded linear functionals
commute with the Bochner integral, so $q(m(X))=\int_Xq\,d\mu$ for every state $q$. Moreover
$\|m(X)\|\le\int_X\|h\|\,d\mu$, $m(X)$ depends only on $X$ up to null sets, and
$m(X\cup Y)=m(X)+m(Y)$ when $X\cap Y$ is null.
\end{definition}

\begin{definition}[Integral ranking]\label{def:integral}
Assume (M1). The \emph{integral ranking} is defined as follows: $X\succeq_\omega Y$ if and only
if the change from $Y$ to $X$
\begin{enumerate}[label=(\roman*)]
  \item is guided at $\omega$, or
  \item is available to no coherent model, and $m_\omega(X)>m_\omega(Y)$.
\end{enumerate}
\end{definition}

\begin{lemma}\label{lem:structure}
Assume (M1), and consider the integral ranking.
\begin{enumerate}[label=(\alph*)]
  \item Whether $X\succeq_\omega Y$ holds depends only on $X\setminus Y$ and $Y\setminus X$ up
  to null sets. In particular $\succeq_\omega$ is a relation on models, and models equal a.e.\
  are tied.
  \item If $X\succeq_\omega Y$ and $X\ne Y$ a.e., then $m_\omega(X)>m_\omega(Y)$. Hence
  distinct models are never tied, and $X\succ_\omega Y$ if and only if $X\succeq_\omega Y$ and
  $X\ne Y$ a.e.
\end{enumerate}
\end{lemma}

\begin{proof}
(a) Guidance, availability and $m_\omega(X)-m_\omega(Y)=m_\omega(X\setminus Y)-m_\omega(Y
\setminus X)$ depend only on these two sets up to null sets. The change between models equal
a.e.\ is null, hence guided at every world.

(b) Let the change add $A$ and remove $B$, with $A\cup B$ non-null. If it is guided at
$\omega$, then
\[
  m_\omega(X)-m_\omega(Y)=\int_A|h(\omega)|\,d\mu(h)+\int_B|h(\omega)|\,d\mu(h),
\]
and the integrand is positive a.e.\ on the non-null set $A\cup B$; so the difference is
positive. Otherwise $m_\omega(X)>m_\omega(Y)$ by clause (ii). If $X\ne Y$ a.e.\ were tied at
$\omega$, then $m_\omega(X)>m_\omega(Y)>m_\omega(X)$. The last claim follows from this and (a).
\end{proof}

\begin{lemma}\label{lem:guidedeverywhere}
If the change that adds $A$ and removes $B$ is guided at every world, then $A\subseteq
D_S$ a.e.\ and $B\subseteq-D_S$ a.e., where $D_S=\{h:h\gneq0\}$ (Lemma~\ref{lem:states}(b)).
Every coherent model contains $A$ a.e.\ and meets $B$ in a null set.
\end{lemma}

\begin{proof}
Let $\Omega_0$ be the countable set of Lemma~\ref{lem:states}(e); it is nonempty, since
$\inf_{\Omega_0}\one=1$. Outside the countable union of the null sets $A\cap\{h(\omega)\le0\}$
and $B\cap\{h(\omega)\ge0\}$, $\omega\in\Omega_0$, every $h\in A$ satisfies $h(\omega)>0$ for
all $\omega\in\Omega_0$, so $h\ne0$ and $\inf_\Omega h=\inf_{\Omega_0}h\ge0$, that is,
$h\gneq0$; likewise every $h\in B$ satisfies $-h\gneq0$. So $A\subseteq D_S$ a.e.\ and
$B\subseteq-D_S$ a.e. Now let $Z$ be a coherent model, and $D$ a coherent set of desirable
gambles with $Z=D$ a.e. By (D2), $D_S\subseteq D$, so $A\setminus Z\subseteq(A\setminus D_S)\cup
(D\setminus Z)$ is null. If $h\in D\cap(-D_S)$, then $-h\in D$ by (D2), and $0=h+(-h)\in D$ by
(D4), contradicting (D1); so $D\cap(-D_S)=\emptyset$, and $B\cap Z\subseteq(B\setminus(-D_S))
\cup(Z\setminus D)$ is null.
\end{proof}

\begin{proposition}\label{prop:axioms}
Assume (M1). The integral ranking satisfies (A1)--(A4), and (A5) for the family of all coherent
models.
\end{proposition}

\begin{proof}
\emph{(A1).} Let $X\ne Y$ a.e., and let the change from $Y$ to $X$ be guided at every world. At
each $\omega$, $X\succeq_\omega Y$ by clause (i), and $X\succ_\omega Y$ by
Lemma~\ref{lem:structure}(b).

\emph{(A2).} By Lemma~\ref{lem:structure}, every comparison $U\succeq_\omega W$ gives
$m_\omega(U)\ge m_\omega(W)$. Suppose $A\succeq_\omega B$, $B\succeq_\omega C$ and $C\succ_\omega
A$. Then $C\ne A$ a.e., since models equal a.e.\ are tied, so $m_\omega(C)>m_\omega(A)\ge
m_\omega(B)\ge m_\omega(C)$, which is impossible.

\emph{(A3).} If $Z\sim_\omega X$, then $Z=X$ a.e.\ by Lemma~\ref{lem:structure}(b), so $Z$ and
$X$ are the same model, and $Y\succeq_\omega Z$ is $Y\succeq_\omega X$.

\emph{(A4).} Let $(G_1,G_2,G_3)$ be a margin triple of order $c$ with margin $\kappa$, let $W$
be a model with $G_1\cup G_2\subseteq W$ a.e.\ and $W\cap G_3$ null, and let $W'$ be the
consolidation. The change from $W$ to $W'$ adds $G_3$ and removes $G_1\cup G_2$, up to null
sets.

\emph{No coherent model could make this change}, by Lemma~\ref{lem:Wincoherent}.

\emph{The score rises at every world.} Fix $\omega$, and put $t\coloneqq\mu(G_1)$,
$\alpha\coloneqq\sup_{x\in G_1}x(\omega)$, $\beta\coloneqq\sup_{y\in G_2}y(\omega)$ and
$\gamma\coloneqq\inf_{z\in G_3}z(\omega)$. These are finite, since the three sets are nonempty
and $z(\omega)\ge(x(\omega)+y(\omega))/c+\kappa$ for all $x\in G_1$, $y\in G_2$, $z\in G_3$; and
$\gamma\ge(\alpha+\beta)/c+\kappa$. Hence
\[
  m_\omega(G_3)\ge\gamma\,ct\ge(\alpha+\beta)t+c\kappa t\ge m_\omega(G_1)+m_\omega(G_2)+c\kappa t,
\]
so $m_\omega(W')-m_\omega(W)=m_\omega(G_3)-m_\omega(G_1)-m_\omega(G_2)>0$. By clause (ii),
$W'\succeq_\omega W$ at every $\omega$.

\emph{(A5).} Let $C$ be a coherent model, and let $Y$ dominate $C$. The change from $C$ to $Y$
removes $C\setminus Y\subseteq C$ and adds $Y\setminus C$, which misses $C$; so the coherent
model $C$ could make it. Since $Y\succeq_\omega C$ at every $\omega$, the change is therefore
guided at every world. By Lemma~\ref{lem:guidedeverywhere}, $C$ contains $Y\setminus C$ a.e., so
$Y\setminus C$ is null; and $C$ meets $C\setminus Y$ in a null set, so $C\setminus Y$ is null.
Thus $Y=C$ a.e., and $Y$ and $C$ are tied at every world by Lemma~\ref{lem:structure}(a),
contradicting dominance.
\end{proof}

\begin{corollary}\label{cor:existence}
For every gamble space $E$ there is a finite Borel measure $\mu$ on $E$ satisfying (M1), (M2),
(M3$'$) and (M4). For every $\mu$ satisfying (M1), the integral ranking satisfies (A1)--(A4),
and (A5) for the family of all coherent models; if $\mu$ also satisfies (M2), its undominated
models are exactly the coherent models.
\end{corollary}

\begin{proof}
The measure of Lemma~\ref{lem:measure} satisfies (M1), (M2), (M3$'$) and (M4). The rest is
Proposition~\ref{prop:axioms} and Corollary~\ref{cor:repair}.
\end{proof}

\section{Coherent dominators under the integral ranking}\label{sec:repair}

Throughout this section, (M1) and (M2) are assumed. By Lemma~\ref{lem:identify}(b), the
coherent models are then exactly the models equal a.e.\ to $D_Q$ for some credal set $Q$.

\begin{definition}[Fractional membership]\label{def:weighted}
For a model $H$ let
\[
  F_H\coloneqq\{f\in L^\infty(\mu):0\le f\le\one_H\text{ a.e.}\},\qquad
  T(f)\coloneqq\int f(h)\,h\,d\mu(h)\in E\quad(f\in L^\infty(\mu)).
\]
Here $\one_H$ is the indicator function of $H$. For a model $Y\subseteq
H$, $\one_Y\in F_H$ and $T(\one_Y)=m(Y)$. For $t\in\R$ write $t^+\coloneqq\max(t,0)$ and
$t^-\coloneqq\max(-t,0)$.
\end{definition}

\begin{lemma}\label{lem:minimax}
Let $H$ be a model and $x\in E$, and put
\[
  v\coloneqq\max_{f\in F_H}\ \min_{q\in S}\ q\bigl(T(f)-x\bigr).
\]
\begin{enumerate}[label=(\alph*)]
  \item Both optima are attained, and $v=\min_{q\in S}\bigl(\int_Hq(h)^+\,d\mu(h)-q(x)\bigr)$,
  where the minimum is also attained.
  \item For each state $q$, the maximisers of $f\mapsto q(T(f))$ over $F_H$ are exactly the
  $f$ equal a.e.\ to $\one_{H\cap\{q>0\}}$.
  \item Let $q^\ast$ attain the minimum in (a), and put $N^\ast\coloneqq H\cap\{q^\ast>0\}$.
  Then every maximiser $f$ in the definition of $v$ equals $\one_{N^\ast}$ a.e., and
  $q(m(N^\ast)-x)\ge v$ for every state $q$.
\end{enumerate}
\end{lemma}

\begin{proof}
\emph{Compactness.} Since $\mu$ is finite, $L^\infty(\mu)$ is the dual of $L^1(\mu)$. The set
$F_H$ is convex, bounded and weak-$\ast$ closed, being cut out by the conditions $\int f\varphi
\,d\mu\ge0$ and $\int(\one_H-f)\varphi\,d\mu\ge0$ for $0\le\varphi\in L^1(\mu)$; so it is
weak-$\ast$ compact.

\emph{Continuity.} The identity map of $E$ lies in $L^1(\mu;E)$, so there are simple functions
$s_k=\sum_ie_{k,i}\one_{E_{k,i}}$ with $\int\|h-s_k(h)\|\,d\mu<1/k$. The maps $T_k(f)\coloneqq
\sum_ie_{k,i}\int_{E_{k,i}}f\,d\mu$ are weak-$\ast$-to-norm continuous, and $\|T(f)-T_k(f)\|
\le1/k$ on $F_H$. So $T$ is weak-$\ast$-to-norm continuous on $F_H$, and $Z\coloneqq T(F_H)$ is
a norm-compact convex subset of $E$.

\emph{Minimax.} Put $\Phi(q,z)\coloneqq q(z-x)$ on $S\times Z$. It is affine in each variable,
weak-$\ast$ continuous in $q$, and norm-continuous in $z$ (Lemma~\ref{lem:states}(a)). Both sets
are compact and convex, so Sion's theorem \cite{Sion} gives $\sup_Z\min_S\Phi=\min_S\sup_Z
\Phi$. The function $z\mapsto\min_{q\in S}\Phi(q,z)=\inf_\Omega(z-x)$ (Lemma~\ref{lem:states}(b))
is $1$-Lipschitz on the compact set $Z$, so it attains its maximum; and $q\mapsto\max_Z\Phi(q,
\cdot)$ is a supremum of continuous functions, hence lower semicontinuous on the compact set
$S$, so it attains its minimum.

\emph{The inner maximum.} For $q\in S$ and $f\in F_H$, $q(T(f))=\int f(h)q(h)\,d\mu\le\int_H
q(h)^+\,d\mu$, with equality if and only if $f=1$ a.e.\ on $H\cap\{q>0\}$ and $f=0$ a.e.\ on
$H\cap\{q<0\}$. Since $\{q=0\}$ is null by (M2), equality holds exactly when $f=\one_{H\cap
\{q>0\}}$ a.e. This proves (b) and the formula in (a).

\emph{The saddle point.} Let $f^\ast$ attain $v$ and $z^\ast\coloneqq T(f^\ast)$. Then
\[
  v=\min_S\Phi(\cdot,z^\ast)\le\Phi(q^\ast,z^\ast)\le\max_Z\Phi(q^\ast,\cdot)=v .
\]
So $f^\ast$ maximises $f\mapsto q^\ast(T(f))$, and $f^\ast=\one_{N^\ast}$ a.e.\ by (b). Hence
$z^\ast=m(N^\ast)$ and $q(m(N^\ast)-x)=\Phi(q,z^\ast)\ge v$ for every state $q$.
\end{proof}

\begin{corollary}\label{cor:consistent}
Let $X$ be a consistent incoherent model, and apply Lemma~\ref{lem:minimax} with $H\coloneqq
\hull(X)$ and $x\coloneqq m(X)$. Then $v>0$, and for any minimiser $q^\ast$ the model
\[
  N_0\coloneqq D_{\conv(Q_X\cup\{q^\ast\})}
\]
is coherent, is equal a.e.\ to $\hull(X)\cap\{q^\ast>0\}$, and satisfies $q(m(N_0)-m(X))\ge v$
for every state $q$.
\end{corollary}

\begin{proof}
By Lemma~\ref{lem:hull}(b), $X\subseteq H$ a.e., so $\one_X\in F_H$ and $v\ge\min_Sq(m(X)-
m(X))=0$. By Lemma~\ref{lem:states}(c), $N_0=D_{Q_X}\cap D_{\{q^\ast\}}=H\cap D_{\{q^\ast\}}$.
This differs from $N^\ast=H\cap\{q^\ast>0\}$ only within $\{q^\ast=0\}$, since $D_{\{q^\ast\}}
\setminus\{q^\ast>0\}\subseteq\{q^\ast=0\}$, and $\{q^\ast=0\}$ is null by (M2); so $N_0=N^\ast$
a.e., and $N_0$ is coherent by Lemma~\ref{lem:DQ}. If $v=0$, then $\one_X$ attains $v$, so
$X=N^\ast=N_0$ a.e.\ by Lemma~\ref{lem:minimax}(c), and $X$ would be coherent. The rest is
Lemma~\ref{lem:minimax}(c), since $m(N_0)=m(N^\ast)$.
\end{proof}

\begin{corollary}\label{cor:losing}
Let $X$ be an inconsistent model.
\begin{enumerate}[label=(\alph*)]
  \item The minimum $v\coloneqq\min_{q\in S}\int_Xq(h)^-\,d\mu(h)$ is attained, and $v>0$.
  \item Let $q^\ast$ attain it, and put $L\coloneqq X\cap\{q^\ast<0\}$ and $X_+\coloneqq
  X\cap\{q^\ast>0\}$. Then $q(m(X_+)-m(X))\ge v$ for every state $q$; $q^\ast\in Q_{X_+}$; and
  $L$ is inconsistent.
\end{enumerate}
\end{corollary}

\begin{proof}
(a) Apply Lemma~\ref{lem:minimax} with $H\coloneqq X$ and $x\coloneqq m(X)$. Since $q(m(X))=
\int_Xq\,d\mu$, the bracket in Lemma~\ref{lem:minimax}(a) is $\int_Xq^-\,d\mu$. It vanishes
exactly when $q\in Q_X$, so it is positive at every state, and its attained minimum is
positive.

(b) Here $N^\ast=X_+$, so $q(m(X_+)-m(X))\ge v$ for every state $q$, by
Lemma~\ref{lem:minimax}(c). The set $X\cap\{q^\ast=0\}$ is null by (M2), so $m(X)=m(X_+)+m(L)$
and $q(m(L))\le-v$ for every state $q$. Also $q^\ast>0$ on $X_+$, so $q^\ast\in Q_{X_+}$. If
$q\in Q_L$, then $q(m(L))=\int_Lq\,d\mu\ge0>-v$, which is impossible; so $L$ is inconsistent.
\end{proof}

\begin{lemma}\label{lem:contraction}
Let $X$ be consistent and $P\coloneqq P_{Q_X}$. If $q\in S$, $a\ge0$ and $q(h)\ge-a\,P(h)$
for a.e.\ $h\in X$, then $q\in Q_X$.
\end{lemma}

\begin{proof}
For $q'\in Q_X$ put $\Psi(q')\coloneqq(q+aq')/(1+a)$, a state. For a.e.\ $h\in X$, $(q+aq')(h)
\ge q(h)+a\,P(h)\ge0$, so $\Psi$ maps $Q_X$ into itself. Its iterates satisfy $\Psi^k(q')-q=
(a/(1+a))^k(q'-q)$, so they converge to $q$ in norm, hence weak-$\ast$. As $Q_X$ is nonempty
and closed, $q\in Q_X$.
\end{proof}

\begin{lemma}\label{lem:trim}
Let $X$ be consistent, $0<\tau<1$ and $Q_\tau\coloneqq(1-\tau)Q_X+\tau S$. For every model
$N\subseteq D_{Q_\tau}$, $Q_{X\setminus N}=Q_X$, and hence $\hull(X\setminus N)=\hull(X)$.
\end{lemma}

\begin{proof}
$Q_X\subseteq Q_{X\setminus N}$ by Lemma~\ref{lem:hull}(a). Let $q\in Q_{X\setminus N}$, put
$a\coloneqq(1-\tau)/\tau$, and write $P\coloneqq P_{Q_X}$. Let $h\in X\cap N$. If $h\in
D^\circ_{Q_\tau}$, then $(1-\tau)P(h)+\tau P_S(h)>0$ by Lemma~\ref{lem:states}(d), so $q(h)\ge
P_S(h)>-a\,P(h)$. Otherwise $h\gneq0$, so $q(h)\ge0$ and $P(h)\ge0$, and $q(h)\ge-a\,P(h)$.
For a.e.\ $h\in X\setminus N$, $q(h)\ge0$ and $P(h)\ge0$ (Lemma~\ref{lem:hull}(b)), so $q(h)\ge-a\,P(h)$.
Lemma~\ref{lem:contraction} gives $q\in Q_X$.
\end{proof}

\begin{lemma}\label{lem:change}
Let $X$ and $N$ be models with $N\subseteq\hull(X\setminus N)$ a.e., and suppose that
$m_\omega(N)>m_\omega(X)$ at every world. Then no coherent model could make the change from $X$
to $N$.
\end{lemma}

\begin{proof}
Put $A\coloneqq N\setminus X$ and $B\coloneqq X\setminus N$, and suppose that a coherent model
could make the change. By Lemma~\ref{lem:identify}(b) it is equal a.e.\ to $D_Q$ for some
credal set $Q$, so $D_Q$ could make the change. Then $B\subseteq D_Q$ a.e., so $Q\subseteq
Q_B$ and $\hull(B)\subseteq D_Q$ by Lemma~\ref{lem:hull}(c). Hence $N\subseteq D_Q$ a.e., so $A\subseteq D_Q$ a.e.; since $A\cap
D_Q$ is null, $A$ is null. Then $m(N)-m(X)=-m(B)$, so $m_\omega(B)<0$ at every world, and $B$
is non-null. Let $q\in Q$. Then $q\ge P_Q\ge0$ a.e.\ on $B$ (Lemma~\ref{lem:DQ}), and
$\{q=0\}$ is null by (M2), so $q>0$ a.e.\ on $B$ and $q(m(B))=\int_Bq\,d\mu>0$. But $m(B)\le0$ pointwise and $q$ is positive, so
$q(m(B))\le0$. This is a contradiction.
\end{proof}

\begin{theorem}\label{thm:repair}
Assume (M1) and (M2). Let $X$ be an incoherent model. Then there are a credal set $Q$ and $\eta>0$ such that
$N\coloneqq D_Q$ satisfies:
\begin{enumerate}[label=(\arabic*)]
  \item $q(m(N)-m(X))\ge\eta$ for every state $q$; in particular $m_\omega(N)\ge
  m_\omega(X)+\eta$ at every world;
  \item $N\subseteq\hull(X\setminus N)$;
  \item no coherent model could make the change from $X$ to $N$.
\end{enumerate}
\end{theorem}

\begin{proof}
\emph{Consistent case.} Put $H\coloneqq\hull(X)=D_{Q_X}$ and $P\coloneqq P_{Q_X}$.
Corollary~\ref{cor:consistent} gives $v>0$, a state $q^\ast$ and the coherent model
$N_0=D_{\conv(Q_X\cup\{q^\ast\})}=H\cap D_{\{q^\ast\}}\subseteq H$ with $q(m(N_0)-m(X))\ge v$
for every state $q$. For $0<\tau<1$ let $Q_\tau\coloneqq(1-\tau)Q_X+\tau S$ and
\[
  N_\tau\coloneqq N_0\cap D_{Q_\tau}=D_{\conv(\conv(Q_X\cup\{q^\ast\})\cup Q_\tau)},
\]
a coherent set of desirable gambles and a model by Lemma~\ref{lem:states}(c),(d) and
Lemma~\ref{lem:DQ}. Since $P\ge P_S$, the function $P_{Q_\tau}=P-\tau(P-P_S)$ does not increase
with $\tau$, so $D_{Q_\tau}$ grows as $\tau$ decreases. Every $h\in N_0\subseteq H$ either
satisfies $h\gneq0$, and then lies in $D_{Q_\tau}$ for every $\tau$, or has $P(h)>0$, and then
lies in $D_{Q_\tau}$ for all small $\tau$. So $N_\tau$ increases to $N_0$ as $\tau$
decreases to $0$, and by (M1) and dominated convergence
\[
  \|m(N_0)-m(N_\tau)\|\le\int_{N_0\setminus N_\tau}\|h\|\,d\mu\to0 .
\]
Fix $\tau$ with $\|m(N_0)-m(N_\tau)\|\le v/2$, and put $N\coloneqq N_\tau$ and $\eta\coloneqq
v/2$. Since states have norm $1$, $q(m(N)-m(X))\ge v-v/2=\eta$ for every state $q$, which is
(1). Lemma~\ref{lem:trim} gives $\hull(X\setminus N)=\hull(X)=H\supseteq N_0\supseteq N$, which
is (2). Lemma~\ref{lem:change} gives (3).

\emph{Inconsistent case.} Corollary~\ref{cor:losing} gives $v>0$, $q^\ast$, $L$ and $X_+$ with
$q(m(X_+)-m(X))\ge v$ for every state $q$, and $q^\ast\in Q_{X_+}$. Put $H\coloneqq\hull(X_+)=
D_{Q_{X_+}}$. Since $q^\ast\in Q_{X_+}$, $P_{Q_{X_+}}\le q^\ast$, so $H\subseteq
\{q^\ast\ge0\}$ (Lemma~\ref{lem:DQ}); and $X_+\subseteq H$ a.e.\ by Lemma~\ref{lem:hull}(b). Apply Lemma~\ref{lem:minimax} with this
$H$ and $x\coloneqq m(X)$, and call its value $v_1$. Since $\one_{X_+}\in F_H$,
\[
  v_1\ge\min_{q\in S}q\bigl(m(X_+)-m(X)\bigr)\ge v>0 .
\]
Let $\pi^\ast$ attain the minimum in Lemma~\ref{lem:minimax}(a), and put
\[
  N\coloneqq D_{\conv(Q_{X_+}\cup\{\pi^\ast\})}=H\cap D_{\{\pi^\ast\}},
\]
a coherent set of desirable gambles and a model (Lemmas~\ref{lem:states}(c) and~\ref{lem:DQ}).
It differs from $H\cap\{\pi^\ast>0\}$ only within $\{\pi^\ast=0\}$, which is null by (M2); so,
by Lemma~\ref{lem:minimax}(c), $q(m(N)-m(X))\ge v_1\eqqcolon\eta$ for every state $q$, which is
(1). Since $N\subseteq H\subseteq\{q^\ast\ge0\}$ and $L\subseteq\{q^\ast<0\}$, we have $L\subseteq X\setminus N$, so $Q_{X\setminus N}\subseteq
Q_L=\emptyset$ by Lemma~\ref{lem:hull}(a) and Corollary~\ref{cor:losing}(b). Hence $X\setminus
N$ is inconsistent and $\hull(X\setminus N)=E\supseteq N$, which is (2). For (3): a coherent
model that could make the change would contain $X\setminus N$ a.e.; by
Lemma~\ref{lem:identify}(b) it is equal a.e.\ to $D_{Q'}$ for a credal set $Q'$, so $X\setminus
N\subseteq D_{Q'}$ a.e., and then $\emptyset\ne Q'\subseteq Q_{X\setminus N}$ by
Lemma~\ref{lem:hull}(c), which is impossible.
\end{proof}

\begin{corollary}\label{cor:repair}
Assume (M1) and (M2). Under the integral ranking, every incoherent model $X$ is strictly
dominated by an undominated coherent model $N$, with $\inf_\omega(m_\omega(N)-m_\omega(X))>0$.
In particular, a model is undominated if and only if it is coherent.
\end{corollary}

\begin{proof}
Let $N=D_Q$ and $\eta$ be as in Theorem~\ref{thm:repair}. Since $N$ is coherent
(Lemma~\ref{lem:DQ}) and $X$ is not,
$N\ne X$ a.e. By Theorem~\ref{thm:repair}(1),(3) and clause (ii) of
Definition~\ref{def:integral}, $N\succeq_\omega X$ at every $\omega$, and
Lemma~\ref{lem:structure}(b) makes each comparison strict. The margin is at least $\eta$. $N$
is undominated by Proposition~\ref{prop:axioms}, as is every coherent model.
\end{proof}

\appendix

\section{Proofs of the setting lemmas}\label{app:setting}

\begin{proof}[Proof of Lemma~\ref{lem:states}]
(a) For $q\in S$ and $h\in E$, $-\|h\|\one\le h\le\|h\|\one$ gives $|q(h)|\le\|h\|$, with
equality at $h=\one$. So $S$ lies in the dual unit ball; it is weak-$\ast$ closed, since
positivity and normalisation are pointwise conditions; so it is weak-$\ast$ compact by the
Banach--Alaoglu theorem. It is metrizable: let $(h_n)_{n\ge1}$ be dense in the unit ball of $E$
and put $d(q,r)\coloneqq\sum_{n\ge1}2^{-n}|q(h_n)-r(h_n)|$ for $q,r\in S$. The series converges
uniformly, so $d$ is weak-$\ast$ continuous on $S\times S$; and $d(q,r)=0$ forces $q=r$ on a
dense subset of the unit ball, hence on $E$. So the identity map from $S$ with the weak-$\ast$
topology to $S$ with the metric $d$ is a continuous bijection from a compact space onto a
Hausdorff space, hence a homeomorphism. Convexity is clear. A credal set is a closed subset of
$S$, hence compact, and $q\mapsto q(h)$ is weak-$\ast$ continuous, so its minimum over a credal
set is attained.

(b) $h-(\inf_\Omega h)\one\ge0$ gives $q(h)\ge\inf_\Omega h$ for every $q\in S$, and
$\inf_{\omega}\delta_\omega(h)=\inf_\Omega h$; as the minimum over $S$ is attained by (a), it
equals $\inf_\Omega h$. Applying this to $-h$ gives $\max_{q\in S}q(h)=-\min_{q\in S}q(-h)=
\sup_\Omega h$. The descriptions of $D_S^\circ$ and $-D_S^\circ$ follow, and $D_S=\{h:h\gneq0\}$
because every $h$ with $\inf_\Omega h>0$ satisfies $h\gneq0$.

(c) Each $q$ has norm $1$, so $|P_Q(h)-P_Q(h')|\le\|h-h'\|$. A minimum of linear functions is
superadditive and positively homogeneous; it is monotone because every state is positive; and
$q(\one)=1$ gives the translation rule. So $D_Q^\circ$ is open and a convex cone, and $C_Q$ is
closed. $D_S^\circ\subseteq D_Q^\circ$ since $P_Q\ge P_S$. If $\sup_\Omega h\le0$, then $q(h)\le0$
for all $q$ by (b), so $h\notin D_Q^\circ$; and $h\gneq0$ would force $h=0$, so $h\notin D_Q$.
The set $\conv(Q_1\cup Q_2)$ is the image of $Q_1\times Q_2\times[0,1]$ under the
continuous map $(q_1,q_2,\lambda)\mapsto(1-\lambda)q_1+\lambda q_2$, hence weak-$\ast$ compact
and convex, and a linear function has the same minimum over $Q_1\cup Q_2$ as over this set; so
$P_{\conv(Q_1\cup Q_2)}=\min(P_{Q_1},P_{Q_2})$. Hence $D^\circ_{\conv(Q_1\cup Q_2)}=
D^\circ_{Q_1}\cap D^\circ_{Q_2}$, and taking the union of each side with $\{g:g\gneq0\}$ gives
$D_{\conv(Q_1\cup Q_2)}=D_{Q_1}\cap D_{Q_2}$.

(d) The set is the continuous image of $Q\times S$ under $(q,s)\mapsto(1-\tau)q+\tau s$, so it
is compact, and it is convex. The minimum over it of $q\mapsto q(h)$ is $(1-\tau)P_Q(h)+\tau
P_S(h)$.

(e) Let $(h_j)$ be dense in $E$. For all $j,l$ choose $\omega_{j,l}$ with $h_j(\omega_{j,l})<
\inf_\Omega h_j+1/l$, and let $\Omega_0$ be the set of all these worlds. If $\|h-h_j\|<\eta$,
then $\inf_{\Omega_0}h\le\inf_{\Omega_0}h_j+\eta=\inf_\Omega h_j+\eta\le\inf_\Omega h+2\eta$;
and $\inf_{\Omega_0}h\ge\inf_\Omega h$ always.
\end{proof}

\begin{proof}[Proof of Lemma~\ref{lem:measure}]
Fix $\omega_0\in\Omega$ and put $E_0\coloneqq\{g\in E:g(\omega_0)=0\}$, a separable closed
subspace. Let $(g_k)$ be dense in $E_0$ and $w_k>0$ with $\sum_kw_k=1$ and $\sum_kw_k\|g_k\|<
\infty$, and put $\nu\coloneqq\sum_k\frac{w_k}2(\delta_{g_k}+\delta_{-g_k})$, a Borel
probability measure on $E_0$ with full support, $\int\|g\|\,d\nu<\infty$ and $\nu(-B)=\nu(B)$
for every Borel $B$. Let $\lambda$ be the standard normal distribution on
$\R$, and let $\mu$ be the image of $\nu\otimes\lambda$ under $\Phi(g,t)\coloneqq g+t\one$. The
map $\Phi:E_0\times\R\to E$ is a linear homeomorphism, with inverse $h\mapsto(h-h(\omega_0)
\one,h(\omega_0))$, and the Borel $\sigma$-algebra of $E_0\times\R$ is the product
$\sigma$-algebra, since both factors are separable metric spaces. So $\mu$ is a Borel
probability measure on $E$.
(M1): $\int\|h\|\,d\mu\le\int\|g\|\,d\nu+\int|t|\,d\lambda<\infty$.
(M3): a nonempty open $W\subseteq E$ has $\Phi^{-1}(W)\supseteq U\times I$ for nonempty open
$U\subseteq E_0$ and $I\subseteq\R$, so $\mu(W)\ge\nu(U)\lambda(I)>0$.
(M2): $P_Q(g+t\one)=P_Q(g)+t$, so $\Phi^{-1}\{P_Q=0\}=\{(g,t):t=-P_Q(g)\}$, and by Fubini's
theorem $\mu\{P_Q=0\}=\int\lambda\{-P_Q(g)\}\,d\nu(g)=0$.
Atomless: for $h\in E$, $\mu\{h\}=\nu\{h-h(\omega_0)\one\}\,\lambda\{h(\omega_0)\}=0$.
Since $\mu$ is atomless, $\mu_c=\mu$, and (M3$'$) is (M3).
Symmetric: $\Phi(-g,-t)=-\Phi(g,t)$, and $\nu\otimes\lambda$ is invariant under $(g,t)\mapsto
(-g,-t)$ because $\nu$ and $\lambda$ are; so $\mu(-B)=\mu(B)$, and (M4) holds.
\end{proof}

\begin{proof}[Proof of Lemma~\ref{lem:tools}]
(a) For each $n\ge1$, fewer than $n\mu(E)+1$ gambles have $\mu\{h\}>1/n$; so $A_\mu$ is
countable, hence Borel, and $\mu_c$ is a finite Borel measure. If $h\in A_\mu$ then
$\mu_c\{h\}=\mu(\emptyset)=0$, and otherwise $\mu_c\{h\}=\mu\{h\}=0$.

(b) Since $B$ is measurable for the completion, there is a Borel set $B_0\subseteq B$ with
$B\setminus B_0$ null. Put $B_1\coloneqq B_0\setminus A_\mu$, a Borel set with $\mu(B_1)=
\mu_c(B)$, and let $\nu(G)\coloneqq\mu(G\cap B_1)$ for Borel $G$. We show that $\nu$ has no
atom in the classical sense: there is no Borel $M$ with $\nu(M)>0$ such that every Borel
$M'\subseteq M$ has $\nu(M')\in\{0,\nu(M)\}$. Suppose $M$ were such a set. Fix $n\ge1$. As $E$
is separable, it is a countable union of closed balls of radius $1/n$, so some such ball $K_n$
has $\nu(M\cap K_n)>0$, hence $\nu(M\cap K_n)=\nu(M)$. The set $M'\coloneqq M\cap\bigcap_nK_n$
satisfies $\nu(M\setminus M')\le\sum_n\nu(M\setminus K_n)=0$, so $\nu(M')=\nu(M)>0$; but $M'$
has diameter at most $2/n$ for every $n$, so it has at most one point, and $\nu$ vanishes on
points by (a). By Sierpi\'nski's theorem \cite{Sierpinski}, $\nu$ takes every value in
$[0,\nu(B_1)]$ on Borel subsets of $B_1$. So there is a Borel $G\subseteq B_1\subseteq
B\setminus A_\mu$ with $\mu(G)=\nu(G)=t$.

(c) If $h\notin\spt X$, there is $\varepsilon>0$ with $\mu_c(X\cap B(h,\varepsilon))=0$; then
every $g\in B(h,\varepsilon/2)$ has $B(g,\varepsilon/2)\subseteq B(h,\varepsilon)$, so
$g\notin\spt X$. Hence $\spt X$ is closed. The open set $E\setminus\spt X$ is covered by open
balls $B$ with $\mu_c(X\cap B)=0$. As a subspace of a separable metric space, it is Lindel\"of,
so countably many of these balls cover it, and $\mu_c(X\setminus\spt X)=0$. If $X\subseteq Y$
a.e., then $\mu_c(X\cap B)\le\mu_c(Y\cap B)$ for every ball $B$, so $\spt X\subseteq\spt Y$.

(d) Let $a\in A_\mu$, and suppose $\inf_\Omega a\le0\le\sup_\Omega a$. By
Lemma~\ref{lem:states}(b) there are states $q_1,q_2$ with $q_1(a)=\inf_\Omega a\le0$ and
$q_2(a)=\sup_\Omega a\ge0$. The function $\lambda\mapsto((1-\lambda)q_1+\lambda q_2)(a)$ is
continuous on $[0,1]$, so some state $q$ on the segment has $q(a)=0$. Then $\{q\}$ is a credal
set, and $\mu\{P_{\{q\}}=0\}\ge\mu\{a\}>0$, contradicting (M2). Hence $\inf_\Omega a>0$ or
$\sup_\Omega a<0$, that is, $a\in D_S^\circ\cup(-D_S^\circ)$ by Lemma~\ref{lem:states}(b).
\end{proof}

\begin{proof}[Proof of Lemma~\ref{lem:identify}]
(a) \emph{Step 1.} $D$ contains $D_S^\circ$, which is open and contains $\one$. So $\inter D$ is a
nonempty open convex cone not containing $0$, and $D\subseteq\cl\inter D$, since a convex set
with nonempty interior lies in the closure of its interior.
\emph{Step 2.} Let $\psi\in E^\ast$ be nonzero and nonnegative on $D$. For $h\ge0$ and $t>0$,
$h+t\one\in D_S^\circ\subseteq D$, so $\psi(h)\ge-t\psi(\one)$; hence $\psi$ is positive. Then
$|\psi(h)|\le\|h\|\psi(\one)$, so $\psi(\one)>0$ and $\psi/\psi(\one)\in K_D$. So the nonzero
functionals nonnegative on $D$ are the positive multiples of members of $K_D$.
\emph{Step 3.} By the Hahn--Banach separation theorem there is $\varphi\in E^\ast$ with
$\varphi>0$ on the open convex set $\inter D$, which misses $0$. By Step 1, $\varphi\ge0$ on $D$,
so $K_D\ne\emptyset$ by Step 2. $K_D$ is weak-$\ast$ closed and convex, so it is a credal set.
\emph{Step 4.} $\cl D$ is a closed convex cone, so by the bipolar theorem and Step 2, $\cl D=
C_{K_D}$. Hence $D\subseteq C_{K_D}$, and $D^\circ_{K_D}$, being open and contained in $\cl D$,
is contained in $\inter\cl D=\inter D\subseteq D$; here $\inter\cl D=\inter D$ because $D$ is
convex with nonempty interior.

(b) By (D3) and (D4), $D$ is a convex cone; by (D1), $0\notin D$; and by (D2) and
Lemma~\ref{lem:states}(b), $D_S^\circ\subseteq D_S\subseteq D$. So (a) applies, and
$D^\circ_{K_D}\subseteq D\subseteq C_{K_D}$. As $\{g:g\gneq0\}\subseteq D$ by (D2), also
$D_{K_D}\subseteq D$. Assume (M2). Then $D\setminus D_{K_D}\subseteq C_{K_D}\setminus
D^\circ_{K_D}=\{P_{K_D}=0\}$, which is null. So $D$ differs from the Borel set $D_{K_D}$
(Lemma~\ref{lem:DQ}) by a subset of a null set; hence $D$ is a model and $D=D_{K_D}$ a.e. If a
model $X$ is coherent, with $X=D$ a.e., then $X=D_{K_D}$ a.e. Conversely, a model equal a.e.\ to
some $D_Q$ is coherent, since $D_Q$ is a coherent set of desirable gambles (Lemma~\ref{lem:DQ}).
\end{proof}

\begin{proof}[Proof of Lemma~\ref{lem:hull}]
(a) Convexity and the inclusion are clear. Since $S$ is metrizable, closedness can be tested
along sequences. If $q_j\in Q_B$ converge to $q$ and $q(h)<0$, then $q_j(h)<0$ for some $j$; so
$B\cap\{q<0\}\subseteq\bigcup_jB\cap\{q_j<0\}$, which is null.

(b) $Q_B$ is compact and metrizable, so it has a dense sequence $(q_j)$. Off a null set, every
$h\in B$ has $q_j(h)\ge0$ for all $j$, hence $q(h)\ge0$ for all $q\in Q_B$, by continuity of
$q\mapsto q(h)$; that is, $P_{Q_B}(h)\ge0$. Under (M2), $\{P_{Q_B}=0\}$ is null, so
$B\subseteq D^\circ_{Q_B}\subseteq D_{Q_B}=\hull(B)$ a.e.

(c) If $B\subseteq D_Q$ a.e., then, as every $q\in Q$ satisfies $q\ge P_Q\ge0$ on $D_Q$
(Lemma~\ref{lem:DQ}), every $q\in Q$ is nonnegative a.e.\ on $B$, so $Q\subseteq Q_B$. Then $B$
is consistent and $P_{Q_B}\le P_Q$, whence $D^\circ_{Q_B}\subseteq D^\circ_Q$ and $\hull(B)=
D_{Q_B}\subseteq D_Q$.
\end{proof}

\begin{thebibliography}{9}

\bibitem{DiestelUhl}
J.~Diestel and J.~J. Uhl.
\newblock \emph{Vector Measures}.
\newblock American Mathematical Society, 1977.

\bibitem{Sierpinski}
W.~Sierpi\'nski.
\newblock Sur les fonctions d'ensemble additives et continues.
\newblock \emph{Fundamenta Mathematicae}, 3:240--246, 1922.

\bibitem{Sion}
M.~Sion.
\newblock On general minimax theorems.
\newblock \emph{Pacific Journal of Mathematics}, 8:171--176, 1958.

\end{thebibliography}

\end{document}
